\documentclass[aspectratio=169,11pt]{beamer}

\usepackage[T1]{fontenc}
\usepackage{lmodern}

\definecolor{Ink}{HTML}{162033}
\definecolor{Muted}{HTML}{657084}
\definecolor{Accent}{HTML}{2364AA}
\definecolor{Paper}{HTML}{F7F8FA}

\setbeamercolor{normal text}{fg=Ink,bg=Paper}
\setbeamercolor{title}{fg=Ink}
\setbeamercolor{subtitle}{fg=Accent}
\setbeamercolor{author}{fg=Muted}
\setbeamercolor{date}{fg=Muted}

\setbeamerfont{title}{size=\fontsize{30}{34}\selectfont,series=\bfseries}
\setbeamerfont{subtitle}{size=\fontsize{16}{19}\selectfont}

\setbeamertemplate{navigation symbols}{}
\setbeamertemplate{headline}{}
\setbeamertemplate{footline}{}

\title{Opening the Black Box of\\Deep Neural Networks via Information}
\subtitle{Section 2: Information Theory of Deep Learning}
\author{Ravid Shwartz-Ziv \and Naftali Tishby}
\date{}

\begin{document}

\begin{frame}{Problem and intuition}
    
\end{frame}

\begin{frame}{Mutual Information and the Information Plane}
    
\end{frame}

\begin{frame}{The Information Bottleneck Objective}
    
\end{frame}

\begin{frame}{Experimental setup}

\begin{itemize}
    \item Synthetic binary classification problem
    \item 12 binary inputs $\Rightarrow 2^{12} = 4096$ possible input patterns
    \item Binary target $Y \in \{0,1\}$ generated probabilistically from $X$
    \item Fully connected DNN, trained with SGD and cross-entropy
    \item No explicit regularization
\end{itemize}

\vspace{0.3cm}

For each hidden layer $T_i$, the authors compute

\[
I(X;T_i)
\qquad \text{and} \qquad
I(T_i;Y)
\]

\[
I(X;T_i) = \text{input information / complexity}
\]

\[
I(T_i;Y) = \text{label information / fitness}
\]

\end{frame}


\begin{frame}{Main result: two phases of training}

\begin{columns}

\begin{column}{0.60\textwidth}
    \centering
    \includegraphics[width=\textwidth]{Presentation1/figures/figure3.png}
\end{column}

\begin{column}{0.40\textwidth}

\textbf{Phase 1: Fitting}
\[
I(T;Y) \uparrow
\]

\begin{itemize}
    \item Label information increases
    \item Fast phase
\end{itemize}

\vspace{0.3cm}

\textbf{Phase 2: Compression}
\[
I(X;T) \downarrow
\]

\begin{itemize}
    \item Input information decreases
    \item Much longer phase
\end{itemize}



Transition from large mean gradient, small variance, to larger fluctuations, smaller mean.
\end{column}
\end{columns}
\end{frame}

\begin{frame}{What does depth change?}

\begin{columns}

\begin{column}{0.62\textwidth}
    \centering
    \includegraphics[width=\textwidth]{Presentation1/figures/figure5.png}
\end{column}

\begin{column}{0.38\textwidth}

\begin{itemize}
    \item More hidden layers $\Rightarrow$ faster convergence
    \item Deeper layers compress faster
    \item Intermediate representations make later compression easier
\end{itemize}

\vspace{0.4cm}

\[
\text{depth}
\rightarrow
\text{more efficient compression}
\]

\end{column}

\end{columns}

\end{frame}

\begin{frame}{Discussion}
    Conclusions, Hypothesis and limitations 
\end{frame}

\end{document}
